\documentclass{article}
\usepackage[T1]{fontenc}
\usepackage{lmodern}
\usepackage{graphicx}
\usepackage{amsmath,amssymb}
\usepackage[a4paper,margin=2cm]{geometry}
\usepackage[absolute,overlay]{textpos}
\setlength{\TPHorizModule}{1mm}
\setlength{\TPVertModule}{1mm}
\usepackage[indonesian]{babel}
\usepackage{hyperref}
\setlength{\emergencystretch}{2em}
% unit: Halaman 51

\begin{document}

% === Halaman 51 (llm) ===

\begin{itemize}
    \item Untuk menyederhanakan TDES, maka langkah di tengah diganti dengan D (mode EDE).
    \item Ada dua versi TDES dengan mode EDE:
    \begin{itemize}
        \item Menggunakan 2 kunci
        \item Menggunaakn 3 kunci
    \end{itemize}
    \item Kenapa EDE? Supaya kompatibel dengan single DES:
\end{itemize}

\end{document}
